% !PRINTED 762
\weberpage{762}
Dieser Hülfssatz gestattet uns zunächst, beim Beweise des Theorems (1) die vereinfachende Voraussetzung zu machen, die Coëfficienten $C_1, C_2, \ldots, C_m$ seien ganze rationale Zahlen.

Angenommen nämlich, es bestehe eine Gleichung

\weq{5}{X_1e^{x_1}+X_2e^{x_2}+\cdots = 0}

mit algebraischen Coëfficienten $X_1, X_2, \ldots$ und von einander verschiedenen Exponenten $x_1, x_2, \ldots$, so können wir immer annehmen, die $X_1, X_2, \ldots$ gehören einem und demselben algebraischen Körper $\Omega$ an.

Wir bezeichnen mit $u_1, u_2, \ldots$ Variable und bilden die Norm der Linearform $X_1u_1+X_2u_2+\cdots$

\weq{6}{N(X_1u_1+X_2u_2+\cdots) = \Phi(u_1,u_2,\ldots),}

die eine ganze homogene Function der Variablen $u$ mit \gesperrt{rationalen Coëfficienten ist}, deren Grad gleich dem Grade des Körpers $\Omega$ ist. Setzen wir in (6)

\[
u_1 = e^{x_1},\ u_2 = e^{x_2},\ \ldots,
\]

so geht jedes Product von Variablen $u$ in eine Grösse von der Form $e^z$ über, worin $z$ eine Summe von Zahlen $x$ ist, und $\Phi(u_1,u_2,\ldots)$ wird ein Ausdruck von der Form $C_1e^{z_1}+C_2e^{z_2}+\cdots$, dessen Coëfficienten rationale Zahlen sind, die, wenn die $X_1, X_2, \ldots$ nicht alle verschwinden, nach 2. auch dann nicht alle Null sein können, wenn die unter einander gleichen unter den Gliedern der Function $\Phi(e^{x_1},e^{x_2},\ldots)$ in ein einziges Glied $Ce^z$ zusammengefasst werden. Wenn nun aber die Gleichung (5) besteht, so ist auch $\Phi(e^{x_1},e^{x_2},\ldots)=0$, und folglich

\[
C_1e^{z_1}+C_2e^{z_2}+\cdots = 0.
\]

Sind unter den $C_1, C_2, \ldots$ gebrochene Zahlen, so können wir sie durch Multiplication der ganzen Gleichung mit dem Hauptnenner in ganze Zahlen verwandeln.

3. \gesperrt{Es braucht also jetzt nur noch bewiesen zu werden, dass die Gleichung (1) für kein System ganzer rationaler Zahlen $C_1, C_2, \ldots, C_m$, die nicht alle verschwinden, bestehen kann.}

Demnach nehmen wir jetzt an, es bestehe eine Gleichung

\weq{7}{A_1e^{x_1}+A_2e^{x_2}+\cdots = 0,}

deren Coëfficienten $A_1, A_2, \ldots$ ganze rationale Zahlen sind,
%CONT>>
